% Figure from MATH 451 notes, section 13. Compile with the notes preamble.
\begin{tikzpicture}[scale=1.3]
  \fill[figB!12] plot[smooth cycle, tension=0.8] coordinates {(0,0) (-1.6,0.5) (-2.3,-0.4) (-1.7,-1.4) (-0.5,-1.1)};
  \draw[figB, thick, dashed] plot[smooth cycle, tension=0.8] coordinates {(0,0) (-1.6,0.5) (-2.3,-0.4) (-1.7,-1.4) (-0.5,-1.1)};
  \node[font=\small, figB] at (-1.5,-0.5) {$C$};
  \foreach \r in {0.9,0.45,0.3} { \draw[black!45, thin, dashed] (0,0) circle (\r); }
  \node[font=\scriptsize, black!60, anchor=west] at (0.64,0.64) {$B(x,1)$};
  \node[font=\scriptsize, black!60, anchor=west] at (0.9,0.05) {$B(x,\tfrac1n)$ all meet $C$};
  \fill[figB] (-0.62,-0.35) circle (1.4pt) node[below left=-1pt, font=\scriptsize, figB] {$s_1$};
  \fill[figB] (-0.33,-0.2) circle (1.4pt) node[below=1pt, font=\scriptsize, figB] {$s_2$};
  \fill[figB] (-0.2,-0.13) circle (1.4pt);
  \fill[figB] (-0.1,-0.07) circle (1.1pt);
  \draw[figR, fill=white, thick] (0,0) circle (1.8pt);
  \node[font=\scriptsize, figR, anchor=south west, fill=white, inner sep=1pt] at (0.06,0.06) {$x \notin C$};
\end{tikzpicture}
